% Source: Solutions/Karan-Gangadharan-PS7-Q1-4.tex (LaTeX source).
\submittedby{Karan Gangadharan}

\begin{proof}
\textbf{(a)} Let $p\colon X\to Y$ be a continuous map. Suppose there exists a continuous map
$f\colon Y\to X$ such that
\[ p\circ f=I_Y, \]
where $I_Y$ is the identity map on $Y$.

We know that $p$ is continuous. Also, since
\[ p(f(y))=y,\qquad\forall y\in Y, \]
the map $p$ is surjective.

Now, let $U\subseteq Y$ be such that $p^{-1}(U)$ is open in $X$. Since $f$ is continuous, we
have $f^{-1}\bigl(p^{-1}(U)\bigr)$ open in $Y$. But
\[ f^{-1}\bigl(p^{-1}(U)\bigr)=(p\circ f)^{-1}(U)=I_Y^{-1}(U)=U. \]
Therefore, $U$ is open in $Y$.

Thus, $p$ is a continuous surjective map satisfying
\[ U\text{ is open in }Y\iff p^{-1}(U)\text{ is open in }X. \]
Hence, $p$ is a quotient map.
\[ \boxed{p\text{ is a quotient map.}} \]

\medskip
\textbf{(b)} Let $A\subseteq X$ and let $r\colon X\to A$ be a retraction such that
\[ r(a)=a,\qquad\forall a\in A. \]
We know that $r$ is continuous by definition. Also, since $r(a)=a$ for every $a\in A$, the map
$r$ is surjective.

Now, let $U\subseteq A$ be such that $r^{-1}(U)$ is open in $X$. Since $A$ has the subspace
topology, we have $r^{-1}(U)\cap A$ open in $A$. But, using the property of retraction,
\begin{align*}
  r^{-1}(U)\cap A &= \{x\in A : r(x)\in U\}\\
                  &= \{x\in A : x\in U\}\\
                  &= U.
\end{align*}
Therefore, $U$ is open in $A$.

Thus, $r$ is a continuous surjective map such that
\[ U\text{ is open in }A\iff r^{-1}(U)\text{ is open in }X. \]
Hence, $r$ is a quotient map.
\[ \boxed{\text{Every retraction is a quotient map.}} \qedhere \]
\end{proof}
